\section{The Hilbert and Quot functors}
\source{fga-explained:page-0117}\source{fga-explained:page-0118}\source{fga-explained:page-0119}\source{fga-explained:page-0120}\source{fga-explained:page-0121}\source{fga-explained:page-0122}\source{fga-explained:page-0123}\source{fga-explained:page-0124}\source{fga-explained:page-0125}\source{fga-explained:page-0126}\source{fga-explained:page-0127}\source{fga-explained:page-0128}\source{fga-explained:page-0129}\source{fga-explained:page-0130}\source{fga-explained:page-0131}\source{fga-explained:page-0132}\source{fga-explained:page-0133}\source{fga-explained:page-0134}\source{fga-explained:page-0135}\source{fga-explained:page-0136}\source{fga-explained:page-0137}\source{fga-explained:page-0138}\source{fga-explained:page-0139}\source{fga-explained:page-0140}\source{fga-explained:page-0141}\source{fga-explained:page-0142}\source{fga-explained:page-0143}\source{fga-explained:page-0144}\source{fga-explained:page-0145}\source{fga-explained:page-0146}\source{fga-explained:page-0147}\source{fga-explained:page-0148}\source{fga-explained:page-0149}\source{fga-explained:page-0150}\source{fga-explained:page-0151}\source{fga-explained:page-0152}

Let $S$ be noetherian and $X\to S$ of finite type.  For an $S$-scheme $T$,
write $X_T=X\times_ST$.  A family of subschemes is a closed subscheme
$Y\subseteq X_T$ flat over $T$; equivalently it is a quotient
$\OO_{X_T}\twoheadrightarrow\OO_Y$ flat over $T$.

\begin{definition}[Hilbert and Quot functors]\label{fga-def-5-1}\source{fga-explained:page-0118}\uses{fga-def-2-1,fga-def-1-7}
The Hilbert functor $\Hilb_{X/S}$ sends $T$ to flat closed subschemes of $X_T$.
For a coherent sheaf $\mathcal E$ on $X$, the Quot functor $\Quot_{\mathcal E/X/S}$
sends $T$ to equivalence classes of quotients $\mathcal E_T\twoheadrightarrow
\mathcal T$ whose target has proper support over $T$ and is flat over $T$;
quotients are equivalent when their kernels coincide.  Pullback makes both
assignments contravariant.
\end{definition}

\begin{remark}\label{fga-rem-5-2}\dcref{5.2}\source{fga-explained:page-0120}\uses{fga-def-5-1,fga-thm-4-23}
Hilbert and Quot functors satisfy fpqc descent and decompose as coproducts over
Hilbert polynomials $P\in\mathbb Q[m]$; the polynomial of a flat family is
locally constant on the base.
\end{remark}

\section{Castelnuovo--Mumford regularity}
\begin{definition}[regularity]\label{fga-def-5-3}\source{fga-explained:page-0124}\uses{}
For a coherent sheaf $\mathcal F$ on $\PP^n_k$, $\mathcal F$ is $m$-regular if
$H^i(\PP^n,\mathcal F(m-i))=0$ for every $i>0$.  Its regularity is
$\reg(\mathcal F)=\min\{m:\mathcal F\text{ is }m\text{-regular}\}$.
\end{definition}\begin{lemma}\label{fga-lem-5-1}\dcref{5.1}\source{fga-explained:page-0125}\uses{fga-def-5-3}
If $\mathcal F$ is $m$-regular on $\PP^n$, then $\mathcal F(m)$ is generated by
global sections, multiplication
$H^0(\mathcal F(m))\otimes H^0(\OO(1))\to H^0(\mathcal F(m+1))$ is
surjective, and $H^i(\mathcal F(r))=0$ for $i>0$ and $r\ge m-i$.
\end{lemma}
\begin{proof}
Choose a hyperplane avoiding the associated points of $\mathcal F$, use the
exact sequence for multiplication by its equation, and induct on $n$.  The
cohomology long exact sequence gives the vanishing and the multiplication
surjectivity; the base case is the module calculation on $\PP^0$.
\end{proof}

\begin{theorem}[uniform regularity]\label{fga-thm-5-3}\dcref{5.3}\source{fga-explained:page-0126}\uses{fga-lem-5-1}
For every Hilbert polynomial $P$ and integers $n\ge0$, there is an integer
$m(P,n)$ such that every coherent sheaf on $\PP^n_k$ with Hilbert polynomial
$P$ is $m(P,n)$-regular, uniformly in the field $k$.
\end{theorem}
\begin{proof}
Induct on $n$.  A hyperplane avoiding associated points gives a short exact
sequence with a sheaf on $\PP^{n-1}$ and a finite-length kernel.  The induction
bound controls the quotient; the Hilbert polynomial controls the length and
therefore a uniform twist killing the kernel.  Apply \cref{fga-lem-5-1}.
\end{proof}

\section{Semicontinuity and base change}
\begin{lemma}\label{fga-lem-5-4}\dcref{5.4}\source{fga-explained:page-0128}\uses{fga-def-5-3}
For a proper morphism $f:X\to S$ and coherent $\mathcal F$, sufficiently large
twists (or sufficiently large terms of a finite complex of flat modules) make
the higher direct images vanish and make formation of $f_*\mathcal F$ commute
with base change.
\end{lemma}
\begin{lemma}\label{fga-lem-5-5}\dcref{5.5}\source{fga-explained:page-0129}\uses{fga-lem-5-4}
For a noetherian base and a coherent sheaf flat over it, the locus where a
cohomology base-change map is an isomorphism is open; on that locus the direct
image is locally free and commutes with arbitrary base change.
\end{lemma}
\begin{theorem}\label{fga-thm-5-6}\dcref{5.6}\source{fga-explained:page-0130}\uses{fga-lem-5-4,fga-lem-5-5}
Let $f:X\to S$ be proper and $\mathcal F$ coherent.  There is a bounded complex
of finite free modules on affine opens of $S$ computing cohomology after every
base change.  Consequently the functions $s\mapsto\dim H^i(X_s,\mathcal F_s)$
are upper semicontinuous.
\end{theorem}
\begin{proof}
Resolve a sufficiently positive twist by sums of $\OO(-r)$, use
\cref{fga-lem-5-4} to replace cohomology by a finite complex of finite modules,
and apply the rank description of cohomology after tensoring with residue
fields.  Rank of a matrix is lower semicontinuous, hence kernel dimensions are
upper semicontinuous.
\end{proof}
\begin{theorem}\label{fga-thm-5-7}\dcref{5.7}\source{fga-explained:page-0130}\uses{fga-thm-5-6}
With $f$ proper and $\mathcal F$ coherent flat over $S$, if the base-change map
$(R^if_*\mathcal F)\otimes k(s)\to H^i(X_s,\mathcal F_s)$ is surjective at $s$,
then it is an isomorphism near $s$ and $R^if_*\mathcal F$ is locally free there.
\end{theorem}
\begin{theorem}\label{fga-thm-5-10}\dcref{5.10}\source{fga-explained:page-0131}\uses{fga-thm-5-7}
For a proper morphism and a coherent sheaf flat over the base, cohomology and
base change commute on the locus where the preceding surjectivity condition
holds; in particular the Hilbert polynomial is locally constant in a flat
family.
\end{theorem}

\section{Generic flatness and flattening stratification}
\begin{theorem}[generic flatness]\label{fga-thm-5-12}\dcref{5.12}\source{fga-explained:page-0133}\uses{fga-def-1-7}
If $S$ is integral noetherian, $X\to S$ is of finite type, and $\mathcal F$ is
coherent on $X$, there is a nonempty open $U\subseteq S$ such that
$\mathcal F|_{X_U}$ is flat over $U$.
\end{theorem}
\begin{proof}
Reduce to $S=\Spec A$, $X=\Spec B$, and a finite $B$-module $M$.  Induct on
the number of generators, use the torsion-free case over the domain $A$, and
invert one nonzero element to make the successive quotients free.  The resulting
localizations give flatness on a nonempty principal open.
\end{proof}
\begin{theorem}[flattening stratification]\label{fga-thm-5-13}\dcref{5.13}\source{fga-explained:page-0134}\uses{fga-thm-5-10,fga-thm-5-12}
For $X\to S$ projective with relatively ample $\OO_X(1)$ and coherent
$\mathcal F$, there is a finite locally closed decomposition
$S=\coprod_P S_P$ such that a map $T\to S$ factors through $S_P$ exactly when
$\mathcal F_T$ is flat and has Hilbert polynomial $P$ on every fiber.
\end{theorem}
\begin{proof}
Choose a common regularity bound, represent the relevant direct images by
locally free modules using base change, and cut out the rank conditions for the
multiplication maps.  Their determinantal locally closed loci have the stated
universal property; induction on twists proves exhaustivity.
\end{proof}

\section{Construction of Quot schemes}
\begin{theorem}[Grassmannian construction]\label{fga-thm-5-14}\dcref{5.14}\source{fga-explained:page-0136}\uses{fga-thm-5-3}
For integers $r>d>0$, the Grassmannian $\Grass(r,d)$ exists as a projective
scheme over $\Spec\ZZ$, represents locally free rank-$d$ quotients of
$\OO^r$, and carries a universal quotient $\OO^r\twoheadrightarrow\mathcal U$.
\end{theorem}
\begin{proof}
Glue the standard affine charts given by invertible $d\times d$ minors of a
matrix of generators.  The transition maps satisfy the cocycle condition;
Pluecker coordinates give a closed embedding into projective space, proving
separatedness and properness.
\end{proof}
\begin{theorem}[Quot construction]\label{fga-thm-5-15}\dcref{5.15}\source{fga-explained:page-0136}\uses{fga-def-5-1,fga-thm-5-3,fga-thm-5-10,fga-thm-5-13,fga-thm-5-14}
Let $S$ be noetherian, $X\to S$ projective, and $\mathcal E$ coherent on $X$.
For every polynomial $P$, the subfunctor $\Quot^P_{\mathcal E/X/S}$ is
represented by a projective $S$-scheme.
\end{theorem}
\begin{proof}
Choose a uniform regularity $m$.  Every quotient with polynomial $P$ is
generated in degree $m$ and is determined by the quotient of the finite locally
free module $H^0(\mathcal E(m))$; hence it gives a point of a Grassmannian.
The kernel and multiplication relations are closed conditions, and the
flattening stratum cuts out exactly polynomial $P$.  The resulting closed
subscheme has the universal quotient and the asserted universal property.
\end{proof}
\begin{theorem}\label{fga-thm-5-16}\dcref{5.16}\source{fga-explained:page-0137}\uses{fga-thm-5-15,fga-def-5-1}
For a projective $X\to S$, the Hilbert functor is represented by the open and
closed union of the Quot schemes for $\mathcal E=\OO_X$; each Hilbert-polynomial
component is projective over $S$.
\end{theorem}
\begin{lemma}\label{fga-lem-5-17}\dcref{5.17}\source{fga-explained:page-0138}\uses{fga-def-5-3,fga-def-5-1}
Tensoring a flat family of quotients by a line bundle preserves flatness and
identifies the corresponding Quot functors; regularity bounds are invariant
under the twist after translating the integer.
\end{lemma}
\begin{theorem}\label{fga-thm-5-18}\dcref{5.18}\source{fga-explained:page-0139}\uses{fga-thm-5-15}
The Quot scheme for a strongly quasi-projective morphism is strongly
quasi-projective; in the projective case it is projective.
\end{theorem}
\begin{lemma}\label{fga-lem-5-19}\dcref{5.19}\source{fga-explained:page-0140}\uses{}
Every coherent sheaf on an open subscheme of a noetherian scheme extends to a
coherent sheaf on the whole scheme after shrinking the ambient open.
\end{lemma}\begin{theorem}\label{fga-thm-5-20}\dcref{5.20}\source{fga-explained:page-0140}\uses{fga-thm-5-18,fga-lem-5-19}
For $X\to S$ quasi-projective and $\mathcal E$ coherent, each
$\Quot^P_{\mathcal E/X/S}$ is represented by a quasi-projective $S$-scheme.
\end{theorem}
\begin{lemma}\label{fga-lem-5-21}\dcref{5.21}\source{fga-explained:page-0141}\uses{fga-def-1-7}
Finite-type flat morphisms over noetherian schemes are open on the flat locus;
flatness is detected by vanishing of the relevant Tor modules and is preserved
under base change.
\end{lemma}
\begin{theorem}\label{fga-thm-5-22}\dcref{5.22}\source{fga-explained:page-0141}\uses{fga-lem-5-21,fga-thm-5-13}
For $f:X\to S$ finite type and $\mathcal F$ coherent, the flattening locus of
$\mathcal F$ is represented by a locally closed subscheme characterized by the
universal flatness property.
\end{theorem}
\begin{theorem}\label{fga-thm-5-23}\dcref{5.23}\source{fga-explained:page-0142}\uses{fga-thm-5-20,fga-thm-5-22}
The Hilbert functor of a quasi-projective $X\to S$ is represented by a
quasi-projective $S$-scheme, with the universal flat family.
\end{theorem}
\begin{lemma}\label{fga-lem-5-24}\dcref{5.24}\source{fga-explained:page-0144}\uses{fga-thm-4-23}
Faithfully flat quasi-compact morphisms are effective for descent of closed
subschemes and of quasi-coherent ideal sheaves.
\end{lemma}
\begin{theorem}\label{fga-thm-5-25}\dcref{5.25}\source{fga-explained:page-0145}\uses{fga-thm-5-16,fga-thm-5-23,fga-lem-5-24}
The Hilbert functor of a projective (respectively quasi-projective) scheme is
represented by a projective (respectively quasi-projective) scheme, and its
universal family is universal under every base change.
\end{theorem}
\begin{theorem}\label{fga-thm-5-26}\dcref{5.26}\source{fga-explained:page-0147}\uses{fga-thm-5-20,fga-lem-5-17}
The same construction applies to relative Quot functors for a morphism that is
strongly quasi-projective over a noetherian base, and is compatible with open
and closed subschemes and with change of the relatively ample line bundle.
\end{theorem}
